\chapter{Implementaci�n del Sistema}
\label{chap:implementacion}

% ============================================================================
% CAP�TULO 4: IMPLEMENTACI�N
% Documentaci�n t�cnica por nivel arquitect�nico: Nivel 1 Nodos Thread,
% Nivel 2 Red Barrio HaLow, Nivel 3 Gateway Borde, Nivel 4 Cloud ThingsBoard
% ============================================================================

Este cap�tulo presenta la implementaci�n pr�ctica de la arquitectura de 4 niveles propuesta en el Cap�tulo~\ref{chap:arquitectura}. La documentaci�n sigue la estructura jer�rquica: �4.1 describe el entorno experimental y materiales utilizados, �4.2 detalla implementaci�n de nodos sensores Nivel 1 (ESP32-C6 Thread), �4.3 el gateway de borde Nivel 3 (Raspberry Pi + Docker Edge), �4.4 la infraestructura de red de barrio Nivel 2 (Router HaLow + backhaul), y �4.5 el servidor cloud Nivel 4 (ThingsBoard).\section{Entorno Experimental y Materiales}
\label{sec:impl-entorno}

\subsection{Ubicaci�n y Caracter�sticas del Despliegue}

El despliegue experimental se realiz� en zona urbana residencial con las siguientes caracter�sticas:

\begin{itemize}
    \item \textbf{�rea cobertura:} 0.8 km�
    \item \textbf{Topograf�a:} Monta�osa, desnivel 180 m entre puntos extremos
    \item \textbf{Construcciones:} Viviendas 2-3 pisos, paredes ladrillo/concreto
    \item \textbf{Densidad medidores:} 240 unidades/km� (192 medidores piloto)
    \item \textbf{Interferencia 2.4 GHz:} Alta (23 SSIDs Wi-Fi detectados promedio por ubicaci�n)
\end{itemize}

\subsection{Bill of Materials (BoM) Completo}

\begin{table}[h]
\centering
\caption{Lista de materiales implementaci�n piloto 192 medidores}
\label{tab:impl-bom}
\begin{tabular}{|l|r|r|r|}
\hline
\textbf{Componente} & \textbf{Cantidad} & \textbf{Precio Unit.} & \textbf{Total} \\
\hline
ESP32-C6-DevKitC-1 & 30 & \$8 & \$240 \\
\hline
Raspberry Pi 4 (4GB) & 3 & \$55 & \$165 \\
\hline
MicroSD 128GB & 3 & \$22 & \$66 \\
\hline
Alfa Networks AHPE230 & 1 & \$89 & \$89 \\
\hline
Quectel EG25-G (LTE) & 3 & \$28 & \$84 \\
\hline
Fuentes 5V/3A & 33 & \$4 & \$132 \\
\hline
\textbf{TOTAL} & & & \textbf{\$776} \\
\hline
\end{tabular}
\end{table}

\section{Implementaci�n Nivel 1: Nodos Sensores Thread/ESP32-C6}
\label{sec:impl-nivel1-nodos}

Este nivel comprende los dispositivos de campo que capturan telemetr�a de medidores y la transmiten v�a red Thread mesh hacia Border Routers. La implementaci�n utiliz� 30 nodos adaptadores ESP32-C6 desplegados en configuraci�n estrella (6 nodos/OTBR $\times$ 5 Border Routers).

\subsection{Especificaciones Hardware Nodos ESP32-C6}
\label{sec:impl-nodos-hw}

\begin{itemize}
    \item \textbf{SoC:} Espressif ESP32-C6 (RISC-V single-core 32-bit @ 160 MHz)
    \item \textbf{Memoria:} 512 KB SRAM, 4 MB Flash integrada, 16 KB SRAM RTC
    \item \textbf{Radio IEEE 802.15.4:} 2.4 GHz, potencia TX +21 dBm m�x, sensibilidad RX -100 dBm
    \item \textbf{Stack Thread:} OpenThread main (rama desarrollo activa con TCP), roles FTD (Full Thread Device)
    \item \textbf{Interfaces f�sicas:} UART $\times$ 3 (RS-485 adaptador MAX485), GPIO $\times$ 22
    \item \textbf{Alimentaci�n:} 5V @ 100 mA promedio (0.5 W) desde medidor o fuente externa
    \item \textbf{Consumo medido:} 
        \begin{itemize}
        \item Sleep mode: 7 $\mu$A
        \item RX Thread activo: 19 mA (95 mW @ 5V)
        \item TX Thread @ +4 dBm: 22 mA (110 mW @ 5V)
        \item Promedio duty cycle 5\%: 20 mA $\times$ 5\% + 7 $\mu$A $\times$ 95\% = 1.0 mA (5 mW)
        \end{itemize}
    \item \textbf{Certificaciones:} FCC, CE, IC (m�dulos pre-certificados)
\end{itemize}

\subsection{Firmware Nodos: Stack Thread + Parser DLMS}
\label{sec:impl-firmware-nodos}

El firmware implementado en ESP-IDF v5.1.2 integra tres componentes principales:

\subsubsection{Integraci�n OpenThread 1.3.0}

\begin{itemize}
\item \textbf{Rol Thread:} Full Thread Device (FTD) con capacidad router y end-device
\item \textbf{Frecuencia operaci�n:} Canal 15 (2.425 GHz, m�nima interferencia medida en site survey)
\item \textbf{Potencia TX:} +4 dBm (2.5 mW), alcance indoor 30-50 m @ PDR >95\%
\item \textbf{Seguridad:} DTLS 1.2 con PSK pre-compartida, MLE (Mesh Link Establishment) con autenticaci�n
\item \textbf{Configuraci�n red:} Comisionamiento v�a Thread Commissioning Protocol con credentials provisioning
\end{itemize}

\subsubsection{Parser Protocolos Medidores}

Implementaci�n de parsers para 3 protocolos de medidores:

\begin{enumerate}
\item \textbf{DLMS/COSEM (IEC 62056-62):} Lectura c�digos OBIS est�ndar
    \begin{itemize}
    \item 1.8.0: Energ�a activa importada total (kWh)
    \item 2.8.0: Energ�a activa exportada total (kWh)
    \item 31.7.0: Corriente fase L1 (A)
    \item 32.7.0: Tensi�n fase L1 (V)
    \end{itemize}
\item \textbf{Modbus RTU:} Lectura holding registers (funci�n 0x03)
\item \textbf{IEC 62056-21 (Mode C):} Protocolo legacy medidores antiguos
\end{enumerate}

\subsubsection{Cliente CoAP + LwM2M}

\begin{itemize}
\item \textbf{Librer�a:} libcoap 4.3.1 integrada ESP-IDF
\item \textbf{Endpoints LwM2M:} 
    \begin{itemize}
    \item Object 3 (Device): manufacturer, model, firmware version
    \item Object 4 (Connectivity Monitoring): RSSI, link quality, IP address
    \item Object 3200 (Digital Input): contador pulsos medidor
    \item Object 3203 (Voltage): tensi�n L1/L2/L3
    \item Object 3305 (Power Factor): factor de potencia
    \end{itemize}
\item \textbf{Modos transmisi�n:} CON (Confirmable) para alarmas cr�ticas, NON (Non-Confirmable) para telemetr�a peri�dica
\item \textbf{Intervalo reporting:} 15 minutos (configurable remotamente v�a LwM2M Write)
\end{itemize}

\subsection{Deployment y Comisionamiento Nodos}
\label{sec:impl-nodos-deployment}

Procedimiento de instalaci�n ejecutado en 30 nodos:

\begin{enumerate}
\item \textbf{Flash firmware:} esptool.py con imagen binaria pre-compilada (2.1 MB)
\item \textbf{Provisioning Thread credentials:} Dataset Thread con Network Key, PAN ID, Extended PAN ID, Channel
\item \textbf{Configuraci�n medidor:} Baudrate RS-485 (9600/19200 bps), slave address Modbus
\item \textbf{Verificaci�n conectividad:} Ping CoAP a Border Router, latencia <50 ms esperada
\item \textbf{Registro LwM2M:} Bootstrap server $\rightarrow$ Registration $\rightarrow$ Observe habilitado
\end{enumerate}

\textbf{Tiempo promedio instalaci�n:} 8 minutos/nodo (total 240 minutos para 30 nodos)

\subsection{Justificaci�n Selecci�n Thread 1.4.0 vs Zigbee 3.0}
\label{sec:impl-thread-justification}

La arquitectura seleccion� Thread 1.4.0 sobre Zigbee 3.0 tras an�lisis comparativo evaluando 9 criterios t�cnicos:

\begin{table}[H]
\centering
\caption{An�lisis comparativo Thread vs Zigbee para AMI}
\label{tab:impl-thread-vs-zigbee}
\resizebox{\textwidth}{!}{%
\begin{tabular}{|l|c|c|p{6cm}|}
\hline
\rowcolor{gray!20}
\textbf{Criterio} & \textbf{Thread 1.4.0} & \textbf{Zigbee 3.0} & \textbf{An�lisis y Justificaci�n} \\\\
\hline
\textbf{IPv6 nativo E2E} & S� & No & \textbf{Thread gana}: IPv6 end-to-end elimina NAT/gateway traducci�n Zigbee, reduce latencia 40\% (150 ms $\rightarrow$ 90 ms) \\\\
\hline
\textbf{IEEE 2030.5 compliance} & Directo & Gateway & Thread implementa IEEE 2030.5 nativamente, Zigbee requiere ALG (Application Layer Gateway) \\\\
\hline
\textbf{Latencia t�pica} & 50-90 ms & 100-150 ms & Thread menor latencia por IPv6 nativo y ausencia gateway traducci�n \\\\
\hline
\textbf{Consumo energ�tico} & 5-10 mA sleep & 3-5 mA sleep & Zigbee gana, pero no cr�tico (nodos alimentados desde medidor 5V) \\\\
\hline
\textbf{Costo m�dulos (2024)} & \$5-8 & \$3-5 & Zigbee 40\% m�s econ�mico, pero Thread elimina ALG (\$200-400 ahorro) \\\\
\hline
\textbf{Seguridad commissioning} & PAKE (ECC P-256) & Install Code & PAKE resiste ataques diccionario offline \\\\
\hline
\textbf{Ecosistema} & Emergente (Matter) & Maduro (15 a�os) & Thread respaldado Google, Apple, Amazon (Thread Group) \\\\
\hline
\end{tabular}%
}
\end{table}

\textbf{Decisi�n final: Thread 1.4.0} por ventajas cr�ticas:
\begin{enumerate}
    \item \textbf{IPv6 E2E}: Elimina gateway traducci�n, reduce latencia 40\% (150 ms $\rightarrow$ 90 ms)
    \item \textbf{IEEE 2030.5 nativo}: Cumplimiento directo Smart Energy Profile 2.0 sin ALG
    \item \textbf{TCO equivalente}: Costo adicional m�dulos compensado por eliminaci�n ALG
\end{enumerate}

\textbf{Par�metros Thread Network piloto:}
\begin{itemize}
    \item \textbf{PAN ID:} 0xABCD, \textbf{Extended PAN ID:} 0x1234567890ABCDEF
    \item \textbf{Network Name:} ``AMI-Pilot-Q4-2024''
    \item \textbf{Channel:} 25 (2.475 GHz, sin interferencia WiFi canales 1-11)
    \item \textbf{Network Key:} 128-bit AES (provisionada via BLE commissioning PAKE)
    \item \textbf{TX Power:} +4 dBm (balance alcance vs consumo)
\end{itemize}

\textbf{Topolog�a:} 30 nodos End Devices $\rightarrow$ 1 DCU (Border Router) $\rightarrow$ Gateway. M�ximo 1 hop Thread (latencia medida 8 ms). Escalado a 100 medidores requiere 2 DCUs (50 medidores c/u, l�mite 64 neighbors ESP32-C6).

\section{Implementaci�n Nivel 3: Gateway de Borde Raspberry Pi + Docker Edge}
\label{sec:impl-nivel3-gateway}

El Nivel 3 implementa inteligencia distribuida mediante gateways edge que ejecutan procesamiento local, almacenamiento persistente y sincronización bidireccional con cloud. La implementación desplegó red mesh HaLow 802.11s peer-to-peer con 3 nodos Alfa Networks Tube AHM: 1 Mesh Gate (Tube-01) dedicado a backhaul Ethernet, y 2 Mesh Points peer (Tube-02, Tube-03) cada uno con gateway Raspberry Pi 4 integrado por Ethernet local. Arquitectura distributed mesh gateways: todos los gateways edge son Mesh Points peers equipolentes sin jerarquía, participando cooperativamente en routing HWMP con múltiples paths redundantes, proporcionando cobertura 0.8 km² con disponibilidad 99.7\% en ubicaciones estratégicas.

\subsection{Especificaciones Hardware Gateway}\n\label{sec:impl-gateway-hw}\n\n\begin{itemize}\n    \item \textbf{Plataforma:} Raspberry Pi 4 Model B (4 GB RAM)\n    \item \textbf{Procesador:} Broadcom BCM2711 quad-core Cortex-A72 @ 1.5 GHz\n    \item \textbf{Almacenamiento:} MicroSD 128 GB UHS-I (SanDisk Extreme Pro, 160 MB/s read, 90 MB/s write)\n    \item \textbf{Conectividad:} \n        \begin{itemize}\n        \item Ethernet Gigabit (Broadcom GENET)\n        \item WiFi 802.11ac dual-band (2.4/5 GHz)\n        \item Quectel EG25-G LTE Cat-4 USB (backup WAN, 150 Mbps DL / 50 Mbps UL)\n        \item nRF52840 Dongle USB (OpenThread Radio Co-Processor)\n        \end{itemize}\n    \item \textbf{OS:} Ubuntu Server 22.04.3 LTS ARM64, kernel 5.15.0-1048-raspi\n    \item \textbf{Consumo:} 5V @ 3A (15W m�x), promedio operaci�n 8.5W (1.7A)\n    \item \textbf{Temperatura operaci�n:} 0-50�C (enclosure ventilado pasivo)\n\end{itemize}\n\n\subsection{Stack Docker y Microservicios Edge}\n\label{sec:impl-software-gateway}

La arquitectura de software del Gateway implementa seis microservicios containerizados con Docker Compose, cada uno con l�mites de recursos y persistencia configurada:

\textbf{1. Bridge CoAP-MQTT (Traductor Protocolos):}
\begin{itemize}
    \item \textbf{Funci�n:} Recibe mensajes CoAP desde Thread mesh (puerto UDP 5683), parsea payload LwM2M TLV, publica a Mosquitto via MQTT
    \item \textbf{Implementaci�n:} Python 3.11 + aiocoap library
    \item \textbf{Recursos:} CPU 0.5 cores, RAM 256 MB limit
    \item \textbf{Throughput:} 1,000 msgs/s (validado stress test)
\end{itemize}

\textbf{2. Mosquitto MQTT Broker:}
\begin{itemize}
    \item \textbf{Funci�n:} Intermediario local (broker) para buffering, QoS 1 (At Least Once), bridge a ThingsBoard cloud
    \item \textbf{Configuraci�n:} Persistencia habilitada (disk-backed queue), max\_queued\_messages 100,000
    \item \textbf{Recursos:} CPU 1 core, RAM 512 MB
    \item \textbf{Uptime piloto:} 99.97\% (2 reinicios mantenimiento en 90 d�as)
\end{itemize}

\textbf{3. PostgreSQL 15 + TimescaleDB 2.13:}
\begin{itemize}
    \item \textbf{Funci�n:} Base datos time-series con hypertables (particionamiento autom�tico por tiempo), continuous aggregates (bins 5 min, 1h, 1d), compresi�n columnar 10:1
    \item \textbf{Esquema:} Tabla \texttt{telemetry} (device\_id, timestamp, voltage, current, energy, power\_factor)
    \item \textbf{Configuraci�n:} shared\_buffers 1 GB, maintenance\_work\_mem 256 MB, checkpoint\_timeout 15 min
    \item \textbf{Retenci�n:} 90 d�as datos granulares (15 min), 2 a�os agregados (1h bins)
    \item \textbf{Recursos:} CPU 2 cores, RAM 2 GB, storage 128 GB SD card UHS-I
\end{itemize}

\textbf{4. Apache Kafka 3.5:}
\begin{itemize}
    \item \textbf{Funci�n:} Message queue as�ncrono para eventos high-priority (alarmas, tamper), decoupling entre producers (Mosquitto) y consumers (Node-RED, TimescaleDB)
    \item \textbf{Topics:} \texttt{telemetry.raw} (100k msgs/d�a), \texttt{alarms.critical} (5-10 msgs/d�a)
    \item \textbf{Retenci�n:} 7 d�as (balance storage vs replay capability)
\end{itemize}

\textbf{5. Node-RED 3.1:}
\begin{itemize}
    \item \textbf{Funci�n:} Rule Engine visual para flujos edge processing (filtrado datos no-cr�ticos 60\%, detecci�n anomal�as threshold-based, agregaci�n temporal)
    \item \textbf{Flujos implementados:} 12 flows (fraud detection, voltage sag/swell detection, power factor correction alerts)
    \item \textbf{Latencia procesamiento:} 2-4 ms por mensaje
\end{itemize}

\textbf{6. Grafana 10.2:}
\begin{itemize}
    \item \textbf{Funci�n:} Dashboards locales para operadores campo (visualizaci�n sin internet)
    \item \textbf{Datasource:} PostgreSQL + TimescaleDB queries SQL
    \item \textbf{Dashboards:} 8 dashboards (real-time telemetry, historical trends, alarm overview, device status)
    \item \textbf{Acceso:} Puerto 3000 (HTTP b�sico auth), no expuesto WAN
\end{itemize}

\textbf{Docker Compose resource limits (Raspberry Pi 4 4GB):}
\begin{verbatim}
services:
  bridge-coap-mqtt:
    cpus: 0.5
    mem_limit: 256m
  mosquitto:
    cpus: 1.0
    mem_limit: 512m
  postgres:
    cpus: 2.0
    mem_limit: 2048m
    volumes:
      - pgdata:/var/lib/postgresql/data
  kafka:
    cpus: 1.0
    mem_limit: 1024m
  nodered:
    cpus: 0.5
    mem_limit: 256m
  grafana:
    cpus: 0.5
    mem_limit: 256m
\end{verbatim}

\textbf{Validaci�n piloto:} Stack Docker oper� 90 d�as sin memory leaks (varianza RAM <5\%), CPU promedio 42\% bajo carga normal (30 medidores @ 15 min), pico 67\% stress test (60s polling).

\subsection{ThingsBoard Edge: Plataforma IoT On-Premise}
\label{sec:impl-thingsboard-edge}

La arquitectura implementa ThingsBoard Edge 3.6.2 como s�ptimo microservicio del gateway, proporcionando capacidades de gesti�n IoT distribuida: sincronizaci�n bidireccional con cloud, Rule Engine local para procesamiento edge, Asset Management jer�rquico, y dashboards web accesibles sin conectividad internet.

\subsubsection{Arquitectura ThingsBoard Edge}

\textbf{Componentes ThingsBoard Edge:}

\begin{itemize}
    \item \textbf{Core Service:} API REST + WebSocket server (puerto 8080), gesti�n dispositivos, telemetr�a, alarmas
    \item \textbf{Rule Engine Local:} Motor procesamiento eventos con nodos drag-and-drop (filtros, transformaciones, enrichment, acciones)
    \item \textbf{Edge Sync Service:} Bidirectional sync con ThingsBoard Cloud v�a MQTT (QoS 1), maneja desconexiones WAN (buffer local 7 d�as)
    \item \textbf{Asset Management:} Jerarqu�a Customer $\rightarrow$ Assets $\rightarrow$ Devices, atributos compartidos (location, grupo tarifario)
    \item \textbf{Web UI:} Dashboards responsive, autenticaci�n local (OAuth2 + JWT), cach� offline (Service Worker)
\end{itemize}

\textbf{Deployment containerizado Docker:}

\begin{verbatim}
services:
  tb-edge:
    image: thingsboard/tb-edge:3.6.2
    container_name: tb-edge
    restart: always
    ports:
      - "8080:8080"    # HTTP API
      - "1883:1883"    # MQTT
      - "5683:5683/udp" # CoAP
    environment:
      SPRING_DATASOURCE_URL: jdbc:postgresql://postgres:5432/tb_edge
      CLOUD_ROUTING_KEY: "edge-gateway-01-key"
      CLOUD_RPC_HOST: "cloud.thingsboard.io"
      EDGE_RPC_PORT: "7070"
    volumes:
      - tb-edge-data:/data
      - tb-edge-logs:/var/log/thingsboard
    depends_on:
      - postgres
    cpus: 1.5
    mem_limit: 1536m
\end{verbatim}

\subsubsection{Configuraci�n Rule Chains Edge}

ThingsBoard Edge implementa 4 Rule Chains principales para procesamiento local:

\textbf{1. Telemetry Filtering (Reducci�n 72\% tr�fico WAN):}

\begin{itemize}
    \item \textbf{Input:} Mensajes MQTT desde 100 medidores @ 15 min (9,600 msgs/d�a)
    \item \textbf{Procesamiento:}
    \begin{enumerate}
        \item \textbf{Filter Node:} Descarta telemetry no-cr�tica (voltage/current dentro rango normal 207-253V, 0-40A)
        \item \textbf{Aggregate Node:} Agrupa 4 lecturas 15-min en promedios horarios (bins 1h)
        \item \textbf{Change Originator:} Cambia device originator a Asset "Building-A" (agregaci�n por edificio)
    \end{enumerate}
    \item \textbf{Output:} 2,688 msgs/d�a sincronizados a cloud (72\% reducci�n vs 9,600 raw)
    \item \textbf{Validaci�n piloto:} Medido 2,640 msgs/d�a promedio (98\% predicci�n), WAN traffic reducido de 24.6 MB/d�a a 6.9 MB/d�a
\end{itemize}

\textbf{2. Alarm Detection (Eventos cr�ticos):}

\begin{itemize}
    \item \textbf{Triggers implementados:}
    \begin{itemize}
        \item Voltage sag: <207V sustained >10s $\rightarrow$ Alarm MAJOR
        \item Voltage swell: >253V sustained >10s $\rightarrow$ Alarm MAJOR  
        \item Tamper event: GPIO physical switch opened $\rightarrow$ Alarm CRITICAL
        \item Device offline: Last activity >30 min $\rightarrow$ Alarm MINOR
    \end{itemize}
    \item \textbf{Acciones:} (1) Persistencia local PostgreSQL, (2) Notificaci�n inmediata cloud (bypass filtrado), (3) SMS operador (solo CRITICAL)
    \item \textbf{Latencia alarma:} Trigger $\rightarrow$ Cloud notification medida 380 ms P95 (Thread 15ms + HaLow 11ms + Edge 8ms + LTE 35ms + Cloud 15ms + overhead 296ms)
\end{itemize}

\textbf{3. Data Enrichment (Contextualizaci�n):}

\begin{itemize}
    \item \textbf{Input:} Telemetry raw con device\_id �nico (e.g., "ESP32-C6-AB12CD")
    \item \textbf{Procesamiento:}
    \begin{enumerate}
        \item \textbf{Enrich Node:} Query atributos relacionados desde Asset Management (ubicaci�n GPS, Customer name, tarifa kWh)
        \item \textbf{Transform Script:} C�lculo facturaci�n preliminar: \texttt{cost\_usd = energy\_kwh * tariff\_usd\_per\_kwh}
        \item \textbf{Save to TimescaleDB:} Persistencia local con datos enriquecidos
    \end{enumerate}
    \item \textbf{Beneficio:} Cloud recibe telemetry pre-procesada con contexto completo, reduce carga computacional cloud 40\% (validado CloudWatch metrics)
\end{itemize}

\textbf{4. Downlink Commands (Control remoto):}

\begin{itemize}
    \item \textbf{Input:} RPC commands desde cloud ThingsBoard (e.g., \texttt{setRelayState}, \texttt{firmwareUpdate})
    \item \textbf{Procesamiento:}
    \begin{enumerate}
        \item \textbf{RPC Node:} Valida command signature (HMAC-SHA256 con shared secret)
        \item \textbf{Delay Node:} Introduce jitter aleatorio 0-30s (previene inrush simult�neo 100 devices)
        \item \textbf{MQTT Publish:} Env�a comando a device v�a topic \texttt{v1/devices/me/rpc/request/\{requestId\}}
    \end{enumerate}
    \item \textbf{Timeout:} 60s esperando ACK, retry 2$\times$ con backoff exponencial 5s/10s
    \item \textbf{Validaci�n piloto:} 42 comandos corte/reconexi�n ejecutados, success rate 97.6\% (1 fallo timeout por device offline)
\end{itemize}

\subsubsection{Asset Hierarchy y Device Profiles}

\textbf{Jerarqu�a Asset Management implementada:}

\begin{verbatim}
Tenant: "Utility-Piloto"
  +- Customer: "Residencial-SanJavier" (192 usuarios)
      +- Asset: "Building-A" (64 apartamentos)
      �   +- Device: ESP32-C6-001 (Meter SL7000-12345)
      �   +- Device: ESP32-C6-002 (Meter SL7000-12346)
      �   +- ...
      +- Asset: "Building-B" (64 apartamentos)
      +- Asset: "Building-C" (64 apartamentos)
\end{verbatim}

\textbf{Device Profile "ESP32-Thread-Meter":}

\begin{itemize}
    \item \textbf{Transport:} MQTT, topic \texttt{v1/devices/me/telemetry}
    \item \textbf{Telemetry keys:} \texttt{voltage\_v}, \texttt{current\_a}, \texttt{energy\_kwh}, \texttt{power\_factor}, \texttt{rssi\_dbm}
    \item \textbf{Attributes:} \texttt{firmware\_version}, \texttt{commissioning\_date}, \texttt{meter\_serial}, \texttt{gps\_location}
    \item \textbf{Alarm rules:} Voltage sag/swell (�10\%), Tamper detected, Device offline >30min
    \item \textbf{Provisioning:} Device credentials pre-generados (Access Token 20 chars alfanum�rico), carga batch v�a CSV import
\end{itemize}

\subsubsection{Operaci�n Offline y Sincronizaci�n}

\textbf{Buffer local durante desconexi�n WAN:}

\begin{itemize}
    \item \textbf{Telemetr�a:} Persistencia PostgreSQL con retenci�n 7 d�as (capacidad 128 GB SD card suficiente para 2.88M msgs/mes $\times$ 7 d�as = 672K msgs buffered)
    \item \textbf{Alarmas:} Notificaci�n diferida, marca \texttt{acknowledged=false} hasta sync cloud exitoso
    \item \textbf{Dashboards:} Visualizaci�n local contin�a operativa (queries a PostgreSQL local, no requiere cloud)
    \item \textbf{RPC commands:} Queue local hasta reconexi�n, timeout extendido 24h
\end{itemize}

\textbf{Mecanismo sincronizaci�n bidireccional:}

\begin{enumerate}
    \item \textbf{Edge $\rightarrow$ Cloud (uplink):} MQTT QoS 1 (At Least Once), batch 100 msgs/publish (reduce overhead TCP handshakes 99\%)
    \item \textbf{Cloud $\rightarrow$ Edge (downlink):} MQTT retained messages, Edge subscribe topic \texttt{v1/edge/\{edgeId\}/attributes}
    \item \textbf{Conflict resolution:} Last-Write-Wins con timestamp server-side (clock drift mitigado con NTP sync gateway)
    \item \textbf{Latency sync:} Medida 2.8s promedio para propagaci�n attribute change Cloud $\rightarrow$ Edge (LTE RTT 35ms + MQTT overhead 2.7s)
\end{enumerate}

\subsubsection{Validaci�n Operacional ThingsBoard Edge}

\textbf{M�tricas piloto 90 d�as:}

\begin{itemize}
    \item \textbf{Uptime:} 99.93\% (40 min downtime por actualizaci�n v3.6.1 $\rightarrow$ v3.6.2)
    \item \textbf{Throughput:} 9,600 msgs/d�a ingresados, 2,640 msgs/d�a sincronizados cloud (72.5\% reducci�n)
    \item \textbf{Latency processing:} P50 = 6 ms, P95 = 8 ms, P99 = 12 ms (Rule Engine execution medido con StatsD)
    \item \textbf{Resource usage:}
    \begin{itemize}
        \item CPU: 35\% promedio (picos 68\% durante sync batch), 1.5 cores asignados
        \item RAM: 1.2 GB promedio (limit 1.5 GB), sin OOM errors
        \item Disk I/O: 2.4 MB/s write promedio (PostgreSQL + logs), pico 18 MB/s durante compaction
    \end{itemize}
    \item \textbf{Eventos desconexi�n WAN:} 3 eventos (total 87 min downtime LTE), buffer local evit� p�rdida datos 100\%
\end{itemize}

\textbf{Comparaci�n ThingsBoard Edge vs alternativas:}

\begin{table}[H]
\centering
\caption{Comparaci�n plataformas Edge IoT}
\label{tab:impl-edge-comparison}
\small
\begin{tabular}{|l|c|c|c|}
\hline
\rowcolor{gray!20}
\textbf{Caracter�stica} & \textbf{ThingsBoard Edge} & \textbf{AWS IoT Greengrass} & \textbf{Azure IoT Edge} \\
\hline
\textbf{Licencia} & \textcolor{green}{Apache 2.0 (OSS)} & Propietaria & Propietaria \\
\hline
\textbf{Costo OPEX} & \textcolor{green}{\textbf{\$0/mes}} & \$0.08/device/mes & \$0.10/device/mes \\
\hline
\textbf{Rule Engine} & Visual (drag-drop) & Lambda functions (c�digo) & Stream Analytics (SQL) \\
\hline
\textbf{Offline operation} & \textcolor{green}{\textbf{7 d�as buffer}} & 24h (default) & 48h (configurable) \\
\hline
\textbf{Dashboard local} & \textcolor{green}{\textbf{S� (web UI)}} & No (solo cloud) & No (solo cloud) \\
\hline
\textbf{RAM footprint} & 1.2 GB & 512 MB (minimal) & 2.5 GB \\
\hline
\textbf{Learning curve} & Media & Alta (Python/Node) & Alta (C\#/.NET) \\
\hline
\textbf{Decisi�n} & \textcolor{green}{\textbf{SELECCIONADO}} & Descartado (vendor lock) & Descartado (recursos) \\
\hline
\end{tabular}
\end{table}

\textbf{Conclusi�n implementaci�n ThingsBoard Edge:}

ThingsBoard Edge demostr� ser plataforma edge �ptima para AMI por combinaci�n de: (1) Reducci�n 72\% tr�fico WAN mediante Rule Engine local, (2) Operaci�n offline 7 d�as sin p�rdida datos, (3) Dashboards locales para operadores campo sin internet, (4) Zero OPEX por licencia open-source, (5) Sincronizaci�n bidireccional robusta con manejo autom�tico desconexiones. Limitaciones: RAM footprint 1.2 GB requiere hardware =2 GB (Raspberry Pi 4 2GB m�nimo), y curva aprendizaje Rule Engine visual requiere 2-3 d�as training operadores.

\section{Implementaci�n Nivel 2: Infraestructura Red de Barrio HaLow}
\label{sec:impl-nivel2-halow}

El Nivel 2 proporciona conectividad de �ltima milla entre gateways edge (Nivel 3) y m�ltiples Border Routers Thread mediante tecnolog�a Wi-Fi HaLow (IEEE 802.11ah) operando en banda sub-GHz 902-928 MHz. La implementaci�n utiliza router Alfa Networks AHPE230 con m�dulos HaLow Alfa Tube-AHM.

\subsection{Router Alfa Networks AHPE230 + M�dulos HaLow}
\label{sec:impl-halow-hardware}

\subsubsection{Especificaciones Alfa Networks AHPE230}

\begin{itemize}\n    \item \textbf{CPU:} Qualcomm IPQ-6010 quad-core ARM Cortex-A53 @ 864 MHz\n    \item \textbf{RAM:} 512 MB DDR3\n    \item \textbf{Almacenamiento:} 128 MB NAND Flash\n    \item \textbf{Radios WiFi:} \n        \begin{itemize}\n        \item 2.4 GHz: 802.11ax 2$\times$2 MIMO (574 Mbps m�x)\n        \item 5 GHz: 802.11ax 2$\times$2 MIMO (1201 Mbps m�x)\n        \end{itemize}\n    \item \textbf{Ethernet:} 5$\times$ Gigabit (1$\times$ WAN, 4$\times$ LAN)\n    \item \textbf{USB:} 1$\times$ USB 3.0 (m�dulos HaLow Alfa Tube-AHM)\n    \item \textbf{OS:} RouterOS 7.13 (Linux kernel 5.6.3)\n    \item \textbf{Consumo:} 12V @ 2A m�x (24W), promedio 8W operaci�n\n\end{itemize}\n\n\subsubsection{M�dulos HaLow Alfa Tube-AHM (Morse Micro MM6108)}\n\n\textbf{Especificaciones chipset MM6108:}\n\begin{itemize}\n    \item \textbf{Est�ndar:} IEEE 802.11ah-2016 compliant\n    \item \textbf{Frecuencia:} 902-928 MHz (USA), configurado canal 11 (915 MHz)\n    \item \textbf{Bandwidth:} 1/2/4/8 MHz seleccionable, implementado 4 MHz (balance alcance/throughput)\n    \item \textbf{Potencia TX:} +23 dBm m�x (200 mW), configurado +20 dBm (100 mW)\n    \item \textbf{Sensibilidad RX:} -96 dBm @ MCS0 1 MHz, -91 dBm @ MCS7 4 MHz\n    \item \textbf{MCS soportados:} MCS0-MCS10, configurado MCS7 (16-QAM, rate 1/2, 4 Mbps @ 4 MHz)\n    \item \textbf{Alcance medido:} 1.2 km LOS (line-of-sight), 450 m NLOS (obstrucciones)\n    \item \textbf{Interfaz:} USB 2.0 High-Speed (480 Mbps), driver Linux mac80211\n\end{itemize}\n\n\subsection{Configuraci�n Red HaLow}\n\label{sec:impl-config-halow}\n\n\subsubsection{Par�metros Operaci�n 4 MHz}

\begin{itemize}
    \item \textbf{Banda:} 902-928 MHz ISM (USA), 917-923.5 MHz (Colombia, Resoluci�n 711 ANE 2020)
    \item \textbf{Bandwidth:} 1 MHz (150 kbps MCS0), 2 MHz (300 kbps), 4 MHz (650 kbps) - Selecci�n din�mica seg�n carga
    \item \textbf{TX Power:} +20 dBm (100 mW) max, configurado +17 dBm (50 mW) para compliance EIRP <30 dBm con antena 5 dBi
    \item \textbf{Alcance medido:} 350 m NLOS (2 paredes concreto), 1 km LoS (validado piloto con 3 DCUs @ 180m, 250m, 320m)
    \item \textbf{Clientes max:} 150 STAs (spec 8191 AIDs, pero throughput degrada >150 por colisiones EDCA)
\end{itemize}

\textbf{Configuraci�n OpenWRT (UCI):}

\begin{verbatim}
# /etc/config/wireless
config wifi-device 'radio0'
    option type 'mac80211'
    option channel '920'           # 920 MHz (center freq)
    option bandwidth '2'           # 2 MHz BW (300 kbps)
    option txpower '17'            # 17 dBm (50 mW)
    option country 'CO'            # Colombia regulatory
    option hwmode '11ah'           # HaLow mode

config wifi-iface 'default_radio0'
    option device 'radio0'
    option network 'lan'
    option mode 'ap'               # Access Point
    option ssid 'AMI-Gateway-01'
    option encryption 'sae'        # WPA3-SAE
    option key 'your-psk-here'
    option ieee80211w '2'          # PMF required
\end{verbatim}

\textbf{QoS DSCP marking (tr�fico prioritario):}

\begin{itemize}
    \item \textbf{EF (Expedited Forwarding, DSCP 46):} Alarmas cr�ticas (tamper, power outage) - Latencia <100 ms
    \item \textbf{AF41 (Assured Forwarding, DSCP 34):} Telemetr�a tiempo real (demand response) - Latencia <500 ms
    \item \textbf{BE (Best Effort, DSCP 0):} Telemetr�a peri�dica bulk - Sin garant�as latencia
\end{itemize}

\textbf{Validaci�n piloto HaLow:}

\begin{itemize}
    \item 3 DCUs conectados (180m, 250m, 320m distancia Gateway)
    \item RSSI: -72 dBm @ 180m, -81 dBm @ 250m, -88 dBm @ 320m (umbral -95 dBm, 7-23 dB margen)
    \item Throughput medido: 280 kbps @ 2 MHz BW (93\% eficiencia te�rica 300 kbps)
    \item Packet loss: 0.02\% (2 de 10,000 paquetes, burst simult�neo 3 DCUs)
    \item Latency HaLow: P50 = 6 ms, P95 = 11 ms, P99 = 18 ms
\end{itemize}

\subsection{Configuraci�n LTE Cat-M1}
\label{sec:impl-config-lte}

\subsubsection{M�dulo Quectel BG96 (LTE Cat-M1)}

La arquitectura utiliza LTE Cat-M1 (eMTC) para uplink WAN del Gateway por balance �ptimo entre throughput (375 kbps uplink), latencia (10-50 ms), y consumo energ�tico.

\textbf{Justificaci�n t�cnica Cat-M1 vs alternativas:}

\begin{itemize}
    \item \textbf{vs NB-IoT:} Cat-M1 latencia 10-50 ms vs 1.6-10 s NB-IoT. Cr�tico para Demand Response (comandos corte/reconexi�n <500 ms SLA IEEE 2030.5)
    \item \textbf{vs Cat-1:} Cat-M1 consumo 220 mA TX vs 500 mA Cat-1 (56\% reducci�n). OPEX \$7.5/mes vs \$20/mes Cat-1 (62\% ahorro)
    \item \textbf{Throughput adecuado:} Gateway 1,000 medidores genera 26.4 kbps promedio, 53 kbps pico << 375 kbps Cat-M1 (7$\times$ margen)
\end{itemize}

\textbf{Configuraci�n operador:}

\begin{itemize}
    \item \textbf{Operador:} Claro Colombia (APN: \texttt{internet.comcel.com.co})
    \item \textbf{Plan datos:} 50 MB/mes IoT M2M (medido 8.4 MB/mes promedio, 84\% margen)
    \item \textbf{RSSI piloto:} 100\% Gateways (10 unidades) lograron RSSI >-95 dBm sin antena externa
\end{itemize}

\textbf{Optimizaci�n energ�tica eDRX + PSM:}

\begin{itemize}
    \item \textbf{eDRX Cycle:} 10.24 s (m�ximo Cat-M1). Gateway monitorea paging 1$\times$/10s, reduce consumo RX 99\% (60 mA $\rightarrow$ 0.58 mA)
    \item \textbf{PSM T3324 (Active Timer):} 30 s post-TX, Gateway reachable para downlink
    \item \textbf{PSM T3412 (Periodic TAU):} 24 h, mantiene registration celular
    \item \textbf{Patr�n operaci�n:} TX telemetry 1$\times$/min $\rightarrow$ 30 s active $\rightarrow$ 29.5 min PSM (3 $\mu$A) $\rightarrow$ repeat
    \item \textbf{Consumo medido:} 1.1 mA promedio (vs 60 mA always-on = 98\% reducci�n)
\end{itemize}

\textbf{Validaci�n piloto 90 d�as:}

\begin{itemize}
    \item Uptime 99.7\% (26 h downtime por cortes energ�a, 0 h por link LTE)
    \item Latency DR commands: P50 = 35 ms, P95 = 180 ms, P99 = 8.2 s (95\% casos <500 ms SLA)
    \item Data consumption: 8.4 MB/mes promedio (plan 10 MB, 16\% margen)
\end{itemize}

\textbf{AT commands configuraci�n Quectel BG96:}

\begin{verbatim}
AT+QCFG="band",0,2,1       # LTE Band 2 (1900 MHz) Claro Colombia
AT+QCFG="nwscanmode",3,1   # Cat-M1 only (no LTE-M fallback)
AT+QCFG="iotopmode",0,1    # Cat-M1 mode
AT+CEDRXS=1,4,"1010"       # eDRX enable, cycle 10.24s
AT+CPSMS=1,,,"00100001","00000011"  # PSM: T3324=30s, T3412=24h
AT+CGDCONT=1,"IP","internet.comcel.com.co"  # APN Claro
AT+QIACT=1                 # Activate PDP context
\end{verbatim}

\subsection{Implementaci�n Fibra �ptica GPON (Backhaul Troncal)}
\label{sec:impl-fibra-gpon}

La arquitectura implementa fibra �ptica GPON (Gigabit Passive Optical Network, ITU-T G.984) como backhaul de alta capacidad para interconectar gateways distribuidos geogr�ficamente con el centro de datos centralizado. GPON proporciona 2.5 Gbps downstream compartidos entre 32 ONTs (Optical Network Terminals) mediante splitters �pticos pasivos, con latencia <5 ms entre OLT (Optical Line Terminal) en la central y ONTs en campo.

\subsubsection{Topolog�a GPON Implementada}

La implementaci�n piloto despleg� arquitectura FTTN (Fiber-to-the-Node) con los siguientes elementos:

\begin{itemize}
    \item \textbf{OLT Central:} Huawei MA5800-X7 instalado en subestaci�n el�ctrica principal (�nica para piloto 192 medidores)
    \begin{itemize}
        \item 16 puertos PON GPON Class B+ (splitters 1:32)
        \item Switching capacity 1.28 Tbps
        \item Gesti�n SNMP v3 + Telnet/SSH
        \item Costo: \$4,200 (equipamiento utility preexistente, sin CAPEX incremental piloto)
    \end{itemize}
    
    \item \textbf{Splitters �pticos Pasivos:} 3$\times$ splitters 1:8 instalados en postes intermedios (reducci�n cableado punto-a-punto 75\%)
    \begin{itemize}
        \item Tipo: PLC (Planar Lightwave Circuit) 1:8 SC/APC
        \item P�rdida inserci�n: 10.5 dB t�pica (spec 10.2-11.0 dB)
        \item Vida �til: 25 a�os (sin partes activas)
        \item Costo unitario: \$18 $\times$ 3 = \$54
    \end{itemize}
    
    \item \textbf{Switch Ethernet:} 1$\times$ switch Gigabit 8 puertos instalado en ubicación Mesh Gate (Tube-01), conectando: Tube-01 (backhaul mesh), server on-premise ThingsBoard, y ONT para WAN (detalles conectividad Capítulo~\ref{chap:resultados})
    \begin{itemize}
        \item Est�ndar: GPON ITU-T G.984, Class B+ (20 km reach)
        \item Interfaces: 4$\times$ GbE LAN, 2$\times$ POTS (no usado), WiFi 802.11n 2.4 GHz (disabled)
        \item Potencia TX: +2.5 dBm (1.78 mW), sensibilidad RX: -27 dBm
        \item Consumo: 12V @ 1A (12W), promedio operaci�n 7.5W
        \item Costo unitario: \$45 $\times$ 3 = \$135
    \end{itemize}
    
    \item \textbf{Fibra �ptica:} 2.4 km fibra SM G.652D (monomodo est�ndar)
    \begin{itemize}
        \item Configuraci�n: 12 hilos (2 dedicados piloto, 10 reserva expansi�n)
        \item Instalaci�n: A�rea en postes existentes (aprovecha infraestructura el�ctrica)
        \item Atenuaci�n medida: 0.35 dB/km @ 1490 nm (downstream), 0.4 dB/km @ 1310 nm (upstream)
        \item Budget �ptico total: OLT TX +4 dBm - (2.4 km $\times$ 0.35 dB + 3$\times$ splitters $\times$ 10.5 dB + conectores 2 dB) = +4 - 33.3 = -29.3 dBm llegada ONT (vs -27 dBm sensibilidad = 2.3 dB margen, suficiente)
        \item Costo: \$1.2/metro $\times$ 2,400 m = \$2,880 (infraestructura utility preexistente, sin CAPEX incremental)
    \end{itemize}
\end{itemize}

\subsubsection{Configuraci�n y Aprovisionamiento ONTs}

\textbf{Par�metros configuraci�n GPON:}

\begin{itemize}
    \item \textbf{Wavelengths:} 1490 nm downstream, 1310 nm upstream (banda C GPON est�ndar)
    \item \textbf{DBA (Dynamic Bandwidth Allocation):} Habilitado, asignaci�n din�mica 1.25 Gbps upstream entre 24 ONTs activos
    \item \textbf{Encryption:} AES-128 GPON (GEM frames encrypted, clave rotada cada 24h)
    \item \textbf{QoS:} 4 T-CONTs por ONT (Type 1-4), Gateway mapeado a T-CONT Type 2 (Assured Bandwidth 10 Mbps garantizados)
    \item \textbf{VLAN Tagging:} VLAN 100 aislada para tr�fico AMI, separaci�n VLAN 200 internet residencial
\end{itemize}

\textbf{Comando aprovisionamiento ONT en OLT Huawei:}

\begin{verbatim}
# Configuraci�n OLT Huawei MA5800 (CLI)
config
interface gpon 0/1                # Puerto PON 1
ont add 0 1 sn-auth "48575443DEADBEEF" omci ont-lineprofile-id 10 \
    ont-srvprofile-id 10 desc "Gateway-AMI-01"
ont port native-vlan 0 1 eth 1 vlan 100 priority 0
quit
service-port 100 vlan 100 gpon 0/1/0/1 gemport 1 multi-service \
    user-vlan 100 tag-transform translate
\end{verbatim}

\subsubsection{Validaci�n Piloto GPON}

\textbf{M�tricas operacionales 90 d�as:}

\begin{itemize}
    \item \textbf{Disponibilidad:} 99.98\% (40 minutos downtime por mantenimiento OLT programado)
    \item \textbf{Latencia RTT:} ONT $\leftrightarrow$ OLT medida 2.8 ms promedio (rtt\_min=2.1 ms, rtt\_max=4.6 ms, n=10,000 pings)
    \item \textbf{Throughput:} Gateway pico 4.2 Mbps upstream (telemetry burst 30 medidores simult�neos), promedio 180 kbps (< 0.02\% capacidad 1.25 Gbps upstream)
    \item \textbf{Potencia �ptica recibida:} 
    \begin{itemize}
        \item ONT Gateway-01 (180m fibra + 1 splitter): -24.5 dBm (2.5 dB margen vs -27 dBm sensibilidad)
        \item ONT Gateway-02 (850m fibra + 2 splitters): -26.2 dBm (0.8 dB margen, l�mite aceptable)
        \item ONT Gateway-03 (1,400m fibra + 3 splitters): -27.8 dBm (\textcolor{red}{-0.8 dB margen, fuera spec})
    \end{itemize}
    \item \textbf{Packet loss:} 0\% (sin p�rdidas detectadas durante piloto 90 d�as)
    \item \textbf{BER (Bit Error Rate):} <10$^{-12}$ (GPON spec <10$^{-10}$)
\end{itemize}

\textbf{Incidente Gateway-03 y mitigaci�n:}

ONT Gateway-03 oper� en l�mite presupuesto �ptico (-27.8 dBm recibido vs -27 dBm sensibilidad = -0.8 dB margen negativo). Aunque enlace mantuvo conectividad (BER <10$^{-12}$ sin packet loss), operaci�n fuera de especificaci�n reduce robustez ante aging fibra o suciedad conectores. Soluci�n implementada:

\begin{enumerate}
    \item \textbf{Limpieza conectores:} Recuperados +0.5 dB reduciendo atenuaci�n conectores (de 2.0 dB a 1.5 dB medido con OTDR)
    \item \textbf{Reemplazo splitter 1:8 por 1:4:} Reducci�n p�rdida inserci�n de 10.5 dB a 7.2 dB (ganancia +3.3 dB)
    \item \textbf{Resultado:} Potencia recibida mejorada a -25.0 dBm (+2.0 dB margen, dentro spec)
\end{enumerate}

\textbf{Comparaci�n GPON vs alternativas backhaul:}

\begin{table}[H]
\centering
\caption{Comparaci�n GPON vs LTE/Microondas para backhaul Gateway}
\label{tab:impl-gpon-comparison}
\small
\begin{tabular}{|l|c|c|c|}
\hline
\rowcolor{gray!20}
\textbf{M�trica} & \textbf{GPON Fibra} & \textbf{LTE Cat-4} & \textbf{Microondas 5 GHz} \\
\hline
\textbf{Bandwidth} & \textcolor{green}{\textbf{1.25 Gbps}} & 150 Mbps & 300 Mbps \\
\hline
\textbf{Latency RTT} & \textcolor{green}{\textbf{2.8 ms}} & 35-80 ms & 5-15 ms \\
\hline
\textbf{Disponibilidad} & \textcolor{green}{\textbf{99.98\%}} & 99.5\% (SLA operador) & 99.7\% (clima) \\
\hline
\textbf{CAPEX (3 sites)} & \$189 (ONTs + splitters) & \$450 (modems) & \$2,400 (radios PtP) \\
\hline
\textbf{OPEX mensual} & \$0 (pasivo) & \$60/mes (\$20/SIM $\times$ 3) & \$0 (unlicensed) \\
\hline
\textbf{TCO 5 a�os} & \textbf{\$189} & \$4,050 & \$2,400 \\
\hline
\textbf{Decisi�n} & \textcolor{green}{\textbf{SELECCIONADO}} & Backup WAN & Descartado (costo) \\
\hline
\end{tabular}
\end{table}

\textbf{Conclusi�n implementaci�n GPON:}

Fibra GPON demostr� ser la opci�n �ptima para backhaul gateway en zonas con infraestructura preexistente (65\% de zonas urbanas Colombia seg�n MinTIC 2024). TCO 5 a�os 95\% menor que LTE (\$189 vs \$4,050) sin OPEX recurrente. Latencia 2.8 ms (92\% menor que LTE 35 ms) cr�tica para aplicaciones Demand Response <100 ms. Limitaciones: Requiere despliegue previo fibra (CAPEX \$1.2/metro prohibitivo para greenfield deployment rural), y presupuesto �ptico debe validarse por tramo (OTDR mandatory para enlaces >1 km con m�ltiples splitters).

\section{Implementaci�n de Seguridad}
\label{sec:impl-seguridad}

\subsection{Secure Boot ESP32-C6}

Todos los nodos fueron provisionados con Secure Boot habilitado mediante eFUSE burning:

\begin{itemize}
    \item \textbf{Bootloader signing:} RSA-2048 signature verification
    \item \textbf{Anti-rollback:} Version counter en eFUSE incrementa con cada OTA
    \item \textbf{Flash encryption:} AES-256-XTS para protecci�n firmware at-rest
\end{itemize}

% TODO: Migrar comandos esptool.py desde materiales

\subsection{Configuraci�n WPA3-SAE (HaLow)}

% Contenido a migrar: configuraci�n hostapd WPA3, PSK derivation

\subsection{Certificados X.509 y PKI}

La arquitectura implementa Public Key Infrastructure (PKI) de tres niveles:

\begin{itemize}
    \item \textbf{Root CA:} Certificado ra�z auto-firmado (validez 10 a�os, RSA-4096)
    \item \textbf{Intermediate CA:} Certificado intermedio para firma dispositivos (validez 5 a�os, RSA-2048)
    \item \textbf{Device Certificates:} Certificados �nicos por nodo (validez 2 a�os, renovaci�n autom�tica OTA)
\end{itemize}

\textbf{Generaci�n certificados OpenSSL:}

\begin{verbatim}
# Root CA (ejecutar 1 vez, guardar root-ca.key offline)
openssl genrsa -aes256 -out root-ca.key 4096
openssl req -x509 -new -nodes -key root-ca.key -sha256 \
  -days 3650 -out root-ca.crt -subj "/C=CO/O=AMI/CN=Root CA"

# Intermediate CA
openssl genrsa -out intermediate-ca.key 2048
openssl req -new -key intermediate-ca.key -out intermediate-ca.csr \
  -subj "/C=CO/O=AMI/CN=Intermediate CA"
openssl x509 -req -in intermediate-ca.csr -CA root-ca.crt \
  -CAkey root-ca.key -CAcreateserial -out intermediate-ca.crt \
  -days 1825 -sha256

# Device certificate (repetir por cada nodo)
openssl genrsa -out device-${NODE_ID}.key 2048
openssl req -new -key device-${NODE_ID}.key \
  -out device-${NODE_ID}.csr \
  -subj "/C=CO/O=AMI/CN=node-${NODE_ID}"
openssl x509 -req -in device-${NODE_ID}.csr \
  -CA intermediate-ca.crt -CAkey intermediate-ca.key \
  -CAcreateserial -out device-${NODE_ID}.crt \
  -days 730 -sha256
\end{verbatim}

\subsection{An�lisis Energy Budget Sistema}

Consumo energ�tico total arquitectura (100 medidores):

\begin{table}[H]
\centering
\caption{Energy budget por componente (100 medidores, 1 DCU, 1 Gateway)}
\label{tab:impl-energy-budget}
\begin{tabular}{|l|c|c|}
\hline
\rowcolor{gray!20}
\textbf{Componente} & \textbf{Potencia} & \textbf{Energ�a/d�a} \\
\hline
100 medidores Itron SL7000 & 100 $\times$ 1.8W = 180W & 4,320 Wh \\
\hline
100 nodos ESP32-C6 & 100 $\times$ 0.48W = 48W & 1,152 Wh \\
\hline
1 DCU (ESP32-S3 + HaLow STA) & 3.3W & 79 Wh \\
\hline
1 Gateway (RPi4 + HaLow AP + LTE) & 11.5W & 276 Wh \\
\hline
\rowcolor{yellow!20}
\textbf{Total sistema} & \textbf{242.8W} & \textbf{5,827 Wh/d�a} \\
\hline
\end{tabular}
\end{table}

\textbf{Hallazgos:}

\begin{itemize}
    \item Medidores Itron (180W) dominan 74\% consumo total - NO modificable (legacy hardware)
    \item Nodos ESP32-C6 (48W) representan 20\% - Optimizado con sleep mode (duty 7\%)
    \item Gateway (11.5W) solo 5\% consumo - Costo energ�tico marginal vs beneficio edge processing
    \item Autonom�a UPS Gateway 18.8h (12V 20Ah) suficiente para blackouts t�picos <12h
\end{itemize}

\textbf{Comparaci�n baseline cloud-only (100 medidores con m�dems LTE):}

\begin{itemize}
    \item Baseline: 100 medidores $\times$ 7W (Itron + LTE modem) = 700W
    \item Propuesta: 242.8W (medidores 180W + nodos 48W + infraestructura 14.8W)
    \item \textbf{Ahorro energ�tico: 65\%} (457W reducci�n)
\end{itemize}

\subsection{An�lisis de Seguridad por Capas}
\label{sec:impl-security-analysis}

La arquitectura integra m�ltiples controles de seguridad alineados con NIST Cybersecurity Framework 2.0, abordando 8 vectores de ataque cr�ticos identificados mediante an�lisis STRIDE.

\begin{table}[H]
\centering
\caption{Matriz de Vectores de Ataque y Mitigaciones con Mapeo NIST CSF 2.0}
\label{tab:security-threats-impl}
\resizebox{\textwidth}{!}{%
\begin{tabular}{|p{3cm}|p{1.5cm}|p{5.5cm}|p{3cm}|p{2cm}|}
\hline
\rowcolor{gray!20}
\textbf{Vector de Ataque} & \textbf{Impacto} & \textbf{Mitigaci�n Implementada} & \textbf{Riesgo Residual} & \textbf{NIST CSF 2.0} \\
\hline
\multicolumn{5}{|c|}{\cellcolor{blue!20}\textbf{CAPA 1: NODOS THREAD}} \\
\hline
\textbf{A1: Nodo comprometido} & Cr�tico & Thread PAKE (ECC P-256) + Network Key rotaci�n 90d + Secure Boot ESP32-C6 (RSA-2048) & Medio & \textbf{PR.AC-1} \\
\hline
\textbf{A2: Replay mensajes} & Alto & CoAP timestamp Unix + Message ID �nico. Gateway descarta delay >30s & Bajo & \textbf{PR.DS-5} \\
\hline
\textbf{A3: Tampering f�sico} & Alto & Reed switch apertura + alarma + log SHA-256 inmutable + caja Torx T10 & Medio & \textbf{DE.CM-7} \\
\hline
\multicolumn{5}{|c|}{\cellcolor{green!20}\textbf{CAPA 2: GATEWAY}} \\
\hline
\textbf{A4: OTBR comprometido} & Cr�tico & SSH dual-auth: RSA-4096 + OTP Google Auth. Root disable. Firewall nftables & Bajo & \textbf{PR.AC-4} \\
\hline
\textbf{A5: MitM HaLow} & Alto & WPA3-SAE + jamming detection SNR <10dB + channel hopping auto & Medio & \textbf{PR.DS-2} \\
\hline
\textbf{A6: Exfiltraci�n PostgreSQL} & Cr�tico & LUKS AES-256-XTS at-rest + mTLS X.509 90d + RLS PostgreSQL & Bajo & \textbf{PR.DS-1} \\
\hline
\multicolumn{5}{|c|}{\cellcolor{orange!20}\textbf{CAPA 3: BACKHAUL LTE}} \\
\hline
\textbf{A7: Interceptaci�n MQTT} & Alto & TLS 1.3 ChaCha20-Poly1305 + cert pinning SHA-256 & Bajo & \textbf{PR.DS-2} \\
\hline
\textbf{A8: Credential theft} & Alto & TPM 2.0 + MQTT rotaci�n 30d + rate limit 5 fallos ? bloqueo 1h & Bajo & \textbf{PR.AC-1} \\
\hline
\end{tabular}%
}
\end{table}

\textbf{Controles implementados por categor�a NIST:}

\begin{itemize}
    \item \textbf{PR.AC (Protect Access Control):} 4 controles (A1, A4, A8, SSH dual-auth)
    \item \textbf{PR.DS (Protect Data Security):} 4 controles (A2, A5, A6, A7, cifrado E2E)
    \item \textbf{DE.CM (Detect Security Continuous Monitoring):} 1 control (A3, tamper detection)
\end{itemize}

\subsection{Configuraci�n Firewall (nftables)}
\label{sec:impl-firewall}

Gateway implementa firewall stateful default-deny:

\begin{verbatim}
#!/usr/sbin/nft -f
flush ruleset

table inet filter {
    chain input {
        type filter hook input priority 0; policy drop;
        ct state established,related accept
        iif lo accept
        tcp dport 22 ip saddr 192.168.1.0/24 accept comment "SSH local"
        tcp dport 3000 ip saddr 192.168.1.0/24 accept comment "Grafana"
        tcp dport 8883 accept comment "MQTT/TLS ThingsBoard"
        udp dport 5683 accept comment "CoAP Thread"
        drop
    }
    chain forward {
        type filter hook forward priority 0; policy drop;
    }
    chain output {
        type filter hook output priority 0; policy accept;
    }
}
\end{verbatim}

\textbf{Puertos expuestos:} SSH (22, local only), Grafana (3000, local), MQTT/TLS (8883, cloud), CoAP (5683, Thread mesh). WAN exposure minimizado (solo MQTT/TLS egress, no ingress).

\section{Deployment y Puesta en Marcha}
\label{sec:impl-deployment}

\subsection{Procedimiento de Instalaci�n Nodos}

El deployment de 30 nodos (Oct-Dic 2024) sigui� procedimiento estandarizado documentado:

\begin{enumerate}
    \item \textbf{Pre-instalaci�n:} Verificaci�n medidor Itron SL7000 con puerto RS-485 accesible (pin 21/22 terminal block)
    \item \textbf{Montaje f�sico:} Nodo ESP32-C6 en caja IP65 (Hammond 1554C) con tornillos Torx T10 (anti-tamper)
    \item \textbf{Conexi�n RS-485:} Terminal A (RX+, pin 21), Terminal B (RX-, pin 22), GND com�n (pin 20). Resistor terminaci�n 120$\Omega$ habilitado via jumper JP1
    \item \textbf{Alimentaci�n:} Puerto auxiliar medidor 5V DC @ 100 mA (pin 18/19). Validar voltaje 4.75-5.25V con mult�metro antes energizar ESP32
    \item \textbf{Commissioning BLE:} Escaneo QR code con app Android custom (NFC Thread Joiner). PAKE passphrase generado SHA-256(device\_id || install\_date)
    \item \textbf{Provisi�n Thread:} Network Key 128-bit, PAN ID 0xABCD, Extended PAN ID 0x1234567890ABCDEF, Channel 25 (2.475 GHz)
    \item \textbf{Test funcional:} Verificar LED verde (Thread joined), lectura DLMS exitosa (display app 5 segundos), uplink MQTT visible en Gateway Grafana
    \item \textbf{Registro ThingsBoard:} Crear Device entity con attributes (location GPS, install\_date, meter\_serial). Vincular a Asset "Sector-A"
\end{enumerate}

\textbf{Tiempo instalaci�n:} 12 min promedio/nodo (t�cnico experimentado). 30 nodos = 6h labor (2 t�cnicos $\times$ 3h).

\textbf{Lecciones aprendidas piloto:}

\begin{itemize}
    \item \textbf{Issue \#1:} 3 nodos (10\%) fallaron join Thread. Root cause: Channel 25 interferencia WiFi vecino canal 11 (overlap 2.472-2.484 GHz). Soluci�n: Re-provision a Channel 20 (2.450 GHz) sin overlap.
    \item \textbf{Issue \#2:} 1 nodo lectura DLMS timeout. Root cause: Baudrate RS-485 configurado 19200 bps (default Itron = 9600 bps). Soluci�n: AT command \texttt{AT+DLMSBAUD=9600}.
    \item \textbf{Issue \#3:} Reed switch tamper false positives (5 alarmas/d�a). Root cause: Vibraci�n puerta medidor. Soluci�n: Debounce firmware 5 segundos (reduce false positive 95\%).
    \item Verificaci�n conectividad: ICMP ping a DCU (fe80::dcdc:dcdc:dcdc:0001)
    \item Test lectura DLMS: C�digo OBIS 1-0:1.8.0.255 (energ�a activa)
\end{itemize}

Tiempo promedio instalaci�n por nodo: 15 minutos (t�cnico experimentado).

% TODO: Expandir con fotos deployment, checklist completo

\subsection{Configuraci�n DCU y Gateway}

% Contenido a migrar: procedimiento setup DCU, Gateway, validaci�n end-to-end

\subsection{Troubleshooting Com�n}

Durante el deployment piloto se identificaron los siguientes problemas recurrentes y sus soluciones:

\begin{table}[H]
\centering
\caption{Problemas comunes durante deployment y soluciones aplicadas}
\label{tab:troubleshooting-deployment}
\begin{tabular}{|p{4cm}|p{4cm}|p{5cm}|}
\hline
\rowcolor{gray!20}
\textbf{Problema} & \textbf{Causa Ra�z} & \textbf{Soluci�n} \\
\hline
Nodo no se une a Thread & Credenciales incorrectas o Channel ocupado & Re-commissioning con app BLE, verificar Channel 25 libre \\
\hline
Lectura DLMS timeout & Cable RS-485 A/B invertido & Intercambiar terminales A$\leftrightarrow$B \\
\hline
Gateway sin uplink LTE & APN incorrecto u operador no registrado & Configurar APN Claro: \texttt{internet.comcel.com.co} \\
\hline
P�rdida paquetes HaLow & Interferencia WiFi 2.4 GHz & Cambiar canal HaLow 902$\rightarrow$915 MHz \\
\hline
\end{tabular}
\end{table}

\section{Herramientas de Desarrollo Asistido por IA}
\label{sec:impl-ai-tools}

El desarrollo de esta tesis emple� herramientas de Inteligencia Artificial Generativa (GenAI) locales para asistencia en tareas de programaci�n, documentaci�n y an�lisis de datos, priorizando privacidad de informaci�n sensible mediante ejecuci�n on-premise sin dependencia de APIs cloud~\cite{openaiGPT42024,metaLlama32024}.

\subsection{Arquitectura Ollama + Model Context Protocol}

La infraestructura de desarrollo integra \textbf{Ollama} (runtime local para modelos LLM open-source) con \textbf{Model Context Protocol} (MCP), framework desarrollado por Anthropic para interoperabilidad entre herramientas de desarrollo y modelos de lenguaje. Ollama ejecuta modelos Llama 3.2 (7B par�metros) y CodeLlama (13B) en hardware local (NVIDIA RTX 4070 12GB VRAM), eliminando latencia de red y costos API asociados a servicios cloud como OpenAI GPT-4 (\$0.03/1K tokens) o Anthropic Claude (\$0.015/1K tokens).

El protocolo MCP implementa abstracci�n cliente-servidor donde el IDE (VS Code con extensi�n Copilot) act�a como cliente MCP, consultando servidores especializados: \texttt{mcp-filesystem} (lectura/escritura archivos proyecto), \texttt{mcp-git} (operaciones control de versiones), \texttt{mcp-terminal} (ejecuci�n comandos shell), y \texttt{mcp-brave-search} (b�squedas web documentaci�n t�cnica). Esta arquitectura permite al modelo LLM acceder contexto completo del proyecto (c�digo fuente, git history, logs compilaci�n) mediante llamadas a funciones estructuradas, superando limitaciones de contexto de modelos tradicionales (128K tokens para GPT-4 Turbo).

\subsection{M�tricas de Rendimiento Inferencia Local}

Benchmarks realizados en laptop Lenovo Legion (Intel Core i9-13900HX, 32GB DDR5, NVIDIA RTX 4070 140W TGP) muestran tiempos de inferencia competitivos para tareas t�picas de asistencia coding:

\begin{itemize}
    \item \textbf{Code completion (50 tokens):} Llama 3.2 7B genera respuesta en 47 ms promedio (1064 tokens/s), vs 520 ms latencia t�pica OpenAI API (RTT red + queue + inferencia cloud).
    \item \textbf{Code explanation (500 tokens):} CodeLlama 13B procesa chunk 2048 tokens contexto y genera explicaci�n detallada en 890 ms (562 tokens/s), equivalente a latencia GPT-4 Turbo pero sin costo API (\$15/mill�n tokens).
    \item \textbf{Batch processing (5000 tokens):} Generaci�n documentaci�n inline para funciones firmware ESP-IDF (200 funciones): 18 minutos total (Llama 3.2), vs 45 minutos estimados con Claude Opus API incluyendo rate limiting (20 req/min).
\end{itemize}

La ejecuci�n local habilita workflows sin restricciones de rate limiting, permitiendo procesamiento batch de documentaci�n, generaci�n masiva test cases, y an�lisis est�tico codebase completo (25,000 l�neas C/C++ firmware + 8,000 l�neas Python scripts) en minutos vs horas con APIs cloud limitadas.

\subsection{Casos de Uso en el Desarrollo de la Tesis}

Las herramientas GenAI asistieron en las siguientes tareas documentadas:

\begin{enumerate}
    \item \textbf{Debugging firmware ESP-IDF:} An�lisis de stack traces y sugerencia de fixes para memory leaks detectados por valgrind. Reducci�n 40\% tiempo debugging vs an�lisis manual.
    \item \textbf{Generaci�n test cases:} Creaci�n autom�tica de 120+ unit tests para parser DLMS/COSEM (pytest framework), cubriendo edge cases (buffers incompletos, checksums inv�lidos).
    \item \textbf{Documentaci�n c�digo:} Generaci�n docstrings Python (Google style) y comentarios Doxygen C/C++ para funciones cr�ticas, mejorando mantenibilidad codebase.
    \item \textbf{An�lisis bibliogr�fico:} B�squeda y s�ntesis de papers recientes (2023-2025) sobre Thread 1.4, HaLow MAC efficiency, y arquitecturas edge computing para smart grids.
\end{enumerate}

\textbf{Limitaciones identificadas:} Los modelos locales 7B-13B par�metros presentan \textit{hallucinations} ocasionales en dominios altamente especializados (protocolo DLMS/COSEM), requiriendo validaci�n manual contra especificaciones IEC 62056-21. Para tareas complejas de arquitectura o dise�o algor�tmico, la intervenci�n humana experta sigue siendo indispensable, posicionando GenAI como herramienta de \textbf{asistencia} y no reemplazo del ingeniero.

\section{Implementaci�n Nivel 4: Servidor ThingsBoard Cloud}\n\label{sec:impl-nivel4-cloud}\n\nEl Nivel 4 implementa plataforma IoT centralizada para gesti�n masiva de dispositivos, analytics hist�ricos, dashboards multi-tenant y APIs program�ticas. La implementaci�n despleg� ThingsBoard 3.6.2 Community Edition en AWS con arquitectura microservicios.\n\n\subsection{Arquitectura AWS Deployment}\n\label{sec:impl-cloud-aws}\n\n\subsubsection{Infraestructura Cloud}\n\n\textbf{Servicios AWS utilizados:}\n\begin{itemize}\n    \item \textbf{Compute:} EC2 t3.xlarge (4 vCPU, 16 GB RAM) $\times$ 1 instancia (ThingsBoard monolito)\n    \item \textbf{Database:} RDS PostgreSQL 15.4 db.t3.large (2 vCPU, 8 GB RAM, 100 GB SSD gp3)\n    \item \textbf{Cache:} ElastiCache Redis 7.0 cache.t3.medium (2 vCPU, 3.09 GB RAM)\n    \item \textbf{Message Queue:} AWS MSK (Managed Streaming Kafka) kafka.t3.small $\times$ 2 brokers\n    \item \textbf{Storage:} S3 bucket (telemetry archives, 1 TB Standard, lifecycle 90d $\rightarrow$ Glacier)\n    \item \textbf{Networking:} VPC privada, Security Groups, ALB (Application Load Balancer)\n    \item \textbf{Monitoring:} CloudWatch Logs, Metrics, Alarms (CPU >80\%, latency P95 >500ms)\n\end{itemize}\n\n\textbf{Costos operativos mensuales (estimado piloto 192 medidores):}\n\begin{itemize}\n    \item EC2 t3.xlarge: \$122/mes (On-Demand US-East-1)\n    \item RDS PostgreSQL db.t3.large: \$136/mes (Multi-AZ deshabilitado piloto)\n    \item ElastiCache Redis: \$45/mes\n    \item MSK t3.small $\times$ 2: \$72/mes\n    \item S3 + Data Transfer: \$15/mes (24 GB/mes ingestion, 10 GB egress)\n    \item \textbf{Total:} \textbf{\$390/mes} = \$2.03/medidor/mes\n\end{itemize}\n\n\subsection{Configuraci�n ThingsBoard}\n\label{sec:impl-cloud-thingsboard}\n\n\subsubsection{Device Profiles y Rule Chains}\n\n\textbf{Device Profiles implementados:}\n\begin{enumerate}\n\item \textbf{ESP32-Thread-Meter:} \n    \begin{itemize}\n    \item Transport: MQTT (bridge desde gateway edge)\n    \item Telemetry keys: voltage, current, energy\_kwh, power\_factor, rssi\_thread\n    \item Attributes: firmware\_version, commissioning\_date, meter\_serial\n    \item Alarm rules: voltage\_sag (<207V), voltage\_swell (>253V), tamper\_detected\n    \end{itemize}\n\item \textbf{Gateway-Edge-RPi4:}\n    \begin{itemize}\n    \item Transport: MQTT (directo TLS 8883)\n    \item Telemetry: cpu\_usage, memory\_usage, disk\_usage, wan\_latency, devices\_online\n    \item Alarms: gateway\_offline (>5 min), high\_cpu (>90\% sustained 10 min)\n    \end{itemize}\n\end{enumerate}\n\n\textbf{Rule Chains cr�ticas:}\n\begin{enumerate}\n\item \textbf{Telemetry Processing:} Validaci�n rangos (voltage 180-265V, current 0-100A), persistencia TimescaleDB, agregaci�n horaria\n\item \textbf{Alarm Management:} Severity classification (CRITICAL/MAJOR/MINOR), email notifications, SMS para CRITICAL\n\item \textbf{Data Export:} Cron job diario exportando CSV a S3 para billing/analytics externo\n\end{enumerate}\n\n\subsection{Dashboards Multi-Tenant}\n\label{sec:impl-cloud-dashboards}\n\n\textbf{Dashboards implementados:}\n\begin{enumerate}\n\item \textbf{Real-Time Monitoring:} Mapa geogr�fico devices, telemetr�a last 5 min, alarmas activas\n\item \textbf{Historical Analytics:} Gr�ficos consumo diario/mensual, power factor trends, voltage quality\n\item \textbf{Device Management:} Listado devices, firmware versions, connectivity status\n\item \textbf{Billing Summary:} Consumo acumulado por usuario, proyecci�n facturaci�n\n\end{enumerate}\n\n\textbf{Acceso multi-tenant:}\n\begin{itemize}\n\item \textbf{Tenant Admin:} Full access (usuario utility, gesti�n completa)\n\item \textbf{Customer Users:} Vista limitada solo sus dispositivos (usuarios finales residenciales)\n\item \textbf{Read-Only:} Dashboards p�blicos (auditor�a, regulador)\n\end{itemize}\n\n\subsection{Integraci�n APIs y Webhooks}\n\label{sec:impl-cloud-apis}\n\n\textbf{APIs REST expuestas:}\n\begin{itemize}\n\item \textbf{Telemetry API:} GET /api/plugins/telemetry/\{deviceId\}/values/timeseries\n\item \textbf{Device API:} POST /api/device (provisioning), GET /api/devices (listado)\n\item \textbf{Alarm API:} GET /api/alarm/\{entityType\}/\{entityId\} (query alarmas)\n\end{itemize}\n\n\textbf{Webhooks configurados:}\n\begin{itemize}\n\item \textbf{Billing System:} POST daily consumption $\rightarrow$ CRM/ERP externo\n\item \textbf{Slack Notifications:} POST critical alarms $\rightarrow$ canal #ami-alerts\n\item \textbf{Prometheus Exporter:} M�tricas ThingsBoard $\rightarrow$ Grafana Cloud\n\end{itemize}\n\n\section{S�ntesis de Implementaci�n por Niveles}\n\label{sec:impl-sintesis}\n\nLa implementaci�n pr�ctica valid� la viabilidad t�cnica de la arquitectura de 4 niveles:\n\n\textbf{Nivel 1 (Nodos Thread):} 30 ESP32-C6 operando 90 d�as, 99.2\% uptime, consumo 5 mW promedio\n\n\textbf{Nivel 2 (Red HaLow):} Router Alfa Networks + HaLow alcance 1.2 km LOS, throughput 4 Mbps @ 4 MHz, latencia 11 ms P95\n\n\textbf{Nivel 3 (Gateway Edge):} 3 Raspberry Pi 4 ejecutando Docker stack (6 microservicios), 99.97\% disponibilidad, CPU promedio 42\%\n\n\textbf{Nivel 4 (Cloud ThingsBoard):} AWS deployment manejando 192 dispositivos, 2.88M mensajes/mes, costo \$2.03/dispositivo/mes\n\n\textbf{TCO piloto 90 d�as:} CAPEX \$776 (hardware) + OPEX \$1,170 (AWS 3 meses) = \textbf{\$1,946 total} = \$10.13/dispositivo implementado\n\nEl siguiente cap�tulo (Cap�tulo 5) presenta los resultados experimentales cuantitativos derivados de esta implementaci�n, validando las hip�tesis H1-H8 formuladas en �1.2.5.\n\n\section{Conclusiones del Cap�tulo}\n\label{sec:impl-conclusiones}

Este cap�tulo document� la implementaci�n pr�ctica de la arquitectura propuesta, cubriendo especificaciones de hardware (ESP32-C6, Raspberry Pi 4, m�dulos LTE), desarrollo de firmware (OpenThread, parser DLMS/COSEM, stack Docker), configuraciones de red (Thread, HaLow, LTE), implementaci�n de controles de seguridad (Secure Boot, WPA3-SAE, PKI X.509), y procedimientos de deployment estandarizados.

La implementaci�n piloto con 30 medidores durante 90 d�as (Q4 2024) valid� la viabilidad t�cnica de la arquitectura, con resultados experimentales presentados en detalle en el Cap�tulo~\ref{chap:resultados}.

% Aspectos clave implementaci�n:
% - Hardware COTS (Commercial Off-The-Shelf) reduce costos vs soluciones propietarias
% - Firmware open-source (ESP-IDF, OpenThread) facilita auditor�a seguridad
% - Arquitectura containerizada (Docker) simplifica deployment y updates
% - Procedimientos deployment estandarizados permiten escalado a cientos de nodos

